\documentclass[11pt]{article}
\usepackage[margin=1.1in]{geometry}
\usepackage{amsmath,amssymb,amsthm,mathtools,booktabs,array,tabularx,longtable}
\usepackage[dvipsnames]{xcolor}
\usepackage[colorlinks=true,linkcolor=MidnightBlue,citecolor=MidnightBlue,urlcolor=MidnightBlue]{hyperref}
\usepackage{fancyhdr}
\pagestyle{fancy}\fancyhf{}\renewcommand{\headrulewidth}{0pt}
\fancyfoot[R]{\thepage}
\fancyfoot[L]{\scriptsize\color{gray}Generated by the Transport Lab from approved parts of its store; not a deposited paper. Papers of record: \href{https://zenodo.org/communities/transport-papers}{Zenodo, Transport Papers}}
\fancypagestyle{plain}{\fancyhf{}\fancyfoot[R]{\thepage}\fancyfoot[L]{\scriptsize\color{gray}Generated by the Transport Lab from approved parts of its store; not a deposited paper. Papers of record: \href{https://zenodo.org/communities/transport-papers}{Zenodo, Transport Papers}}\renewcommand{\headrulewidth}{0pt}}
\newtheorem{theorem}{Theorem}[section]
\newtheorem{proposition}[theorem]{Proposition}
\newtheorem{lemma}[theorem]{Lemma}
\newtheorem{corollary}[theorem]{Corollary}
\newtheorem{conjecture}[theorem]{Conjecture}
\newtheorem{refutation}[theorem]{Refutation}
\newtheorem{prediction}[theorem]{Prediction}
\newtheorem{principle}[theorem]{Principle}
\newtheorem{observation}[theorem]{Observation}
\newtheorem{axiom}[theorem]{Axiom}
\theoremstyle{definition}
\newtheorem{definition}[theorem]{Definition}
\newtheorem{fact}[theorem]{Fact}
\newtheorem{assembled}[theorem]{Assembly}
\newtheorem{postulate}[theorem]{Postulate}
\newtheorem{premise}[theorem]{Premise}
\newtheorem{hypothesis}[theorem]{Hypothesis}
\newtheorem{remark}[theorem]{Remark}
\newtheorem{example}[theorem]{Example}
\theoremstyle{plain}
\setlength{\marginparwidth}{0.8in}
\newcommand{\changedmark}[1]{\marginpar{\raggedright\scriptsize\color{BrickRed}#1}}
\providecommand{\Aa}{}\renewcommand{\Aa}{\mathcal{A}}
\providecommand{\Ad}{}\renewcommand{\Ad}{\widetilde{\mathrm{Ad}}}
\providecommand{\Aev}{}\renewcommand{\Aev}{\mathcal{A}^{+}}
\providecommand{\Aod}{}\renewcommand{\Aod}{\mathcal{A}^{-}}
\providecommand{\B}{}\renewcommand{\B}{\mathrm{B}}
\providecommand{\C}{}\renewcommand{\C}{\mathbb{C}}
\providecommand{\CAR}{}\renewcommand{\CAR}{\mathrm{CAR}}
\providecommand{\CS}{}\renewcommand{\CS}{{\mathbb{S}^{+}}}
\providecommand{\CSm}{}\renewcommand{\CSm}{{\mathbb{S}^{-}}}
\providecommand{\Cl}{}\renewcommand{\Cl}{\mathrm{Cl}}
\providecommand{\Con}{}\renewcommand{\Con}{\mathcal{C}}
\providecommand{\Dl}{}\renewcommand{\Dl}{\Delta}
\providecommand{\EN}{}\renewcommand{\EN}{E_{N}}
\providecommand{\End}{}\renewcommand{\End}{\mathrm{End}}
\providecommand{\F}{}\renewcommand{\F}{\mathbb{F}}
\providecommand{\G}{}\renewcommand{\G}{\Gamma}
\providecommand{\HH}{}\renewcommand{\HH}{\mathbb{H}}
\providecommand{\IS}{}\renewcommand{\IS}{I\sigma_3}
\providecommand{\J}{}\renewcommand{\J}{\mathbf{J}}
\providecommand{\LH}{}\renewcommand{\LH}{L_{H}}
\providecommand{\Mat}{}\renewcommand{\Mat}{\mathrm{M}}
\providecommand{\Mo}{}\renewcommand{\Mo}{M^{\circ}}
\providecommand{\Mp}{}\renewcommand{\Mp}{M^{+}}
\providecommand{\Msun}{}\renewcommand{\Msun}{M_{\odot}}
\providecommand{\Mt}{}\renewcommand{\Mt}{\widetilde{M}}
\providecommand{\Ncone}{}\renewcommand{\Ncone}{\mathcal{N}}
\providecommand{\Neff}{}\renewcommand{\Neff}{N_{\mathrm{eff}}}
\providecommand{\PCS}{}\renewcommand{\PCS}{\mathbb{P}{\mathbb{S}^{+}}}
\providecommand{\PT}{}\renewcommand{\PT}{\mathbb{PT}}
\providecommand{\Pin}{}\renewcommand{\Pin}{\mathrm{Pin}}
\providecommand{\Q}{}\renewcommand{\Q}{\mathbb{Q}}
\providecommand{\R}{}\renewcommand{\R}{\mathbb{R}}
\providecommand{\Ric}{}\renewcommand{\Ric}{\mathrm{Ric}}
\providecommand{\SBH}{}\renewcommand{\SBH}{S_{\mathrm{BH}}}
\providecommand{\Sc}{}\renewcommand{\Sc}{\mathbb{S}}
\providecommand{\Sch}{}\renewcommand{\Sch}{\mathrm{P}}
\providecommand{\Sig}{}\renewcommand{\Sig}{\Sigma}
\providecommand{\Som}{}\renewcommand{\Som}{S_\omega}
\providecommand{\Sp}{}\renewcommand{\Sp}{\mathbb{S}}
\providecommand{\Sph}{}\renewcommand{\Sph}{\mathrm{S}}
\providecommand{\Spi}{}\renewcommand{\Spi}{S_\pi}
\providecommand{\Spin}{}\renewcommand{\Spin}{\mathrm{Spin}}
\providecommand{\Sth}{}\renewcommand{\Sth}{S_{\mathrm{th}}}
\providecommand{\T}{}\renewcommand{\T}{\mathbb{T}}
\providecommand{\TH}{}\renewcommand{\TH}{T_{\mathrm{H}}}
\providecommand{\Tr}{}\renewcommand{\Tr}{\operatorname{Tr}}
\providecommand{\UHF}{}\renewcommand{\UHF}{\mathrm{UHF}}
\providecommand{\Z}{}\renewcommand{\Z}{\mathbb{Z}}
\providecommand{\Zc}{}\renewcommand{\Zc}{\mathrm{Z}}
\providecommand{\cmg}{}\renewcommand{\cmg}{\mathrm{cm^{2}\,g^{-1}}}
\providecommand{\dS}{}\renewcommand{\dS}{\mathrm{dS}}
\providecommand{\dsc}{}\renewcommand{\dsc}{\mathrm{dsc}}
\providecommand{\dtPG}{}\renewcommand{\dtPG}{dt}
\providecommand{\dtS}{}\renewcommand{\dtS}{dt_{\mathrm{S}}}
\providecommand{\ellv}{}\renewcommand{\ellv}{\ell}
\providecommand{\er}{}\renewcommand{\er}{e_r}
\providecommand{\et}{}\renewcommand{\et}{e_t}
\providecommand{\ev}{}\renewcommand{\ev}{\mathrm{ev}}
\providecommand{\gl}{}\renewcommand{\gl}{\mathfrak{gl}}
\providecommand{\hatV}{}\renewcommand{\hatV}{\widehat{V}}
\providecommand{\hb}{}\renewcommand{\hb}{\bar{h}}
\providecommand{\hf}{}\renewcommand{\hf}{\hat{f}}
\providecommand{\id}{}\renewcommand{\id}{\mathrm{id}}
\providecommand{\kpc}{}\renewcommand{\kpc}{\mathrm{kpc}}
\providecommand{\lP}{}\renewcommand{\lP}{\ell_{P}}
\providecommand{\mD}{}\renewcommand{\mD}{\mathbf{m}_{D}}
\providecommand{\mb}{}\renewcommand{\mb}{\bar m}
\providecommand{\mnu}{}\renewcommand{\mnu}{\Sigma m_{\nu}}
\providecommand{\nInf}{}\renewcommand{\nInf}{n_\infty}
\providecommand{\nO}{}\renewcommand{\nO}{n_0}
\providecommand{\ninf}{}\renewcommand{\ninf}{n_{\infty}}
\providecommand{\nz}{}\renewcommand{\nz}{n_{0}}
\providecommand{\proj}{}\renewcommand{\proj}{\operatorname{proj}}
\providecommand{\pv}{}\renewcommand{\pv}{\pi}
\providecommand{\rad}{}\renewcommand{\rad}{\mathrm{rad}}
\providecommand{\sgn}{}\renewcommand{\sgn}{\operatorname{sgn}}
\providecommand{\shA}{}\renewcommand{\shA}{\mathcal{A}}
\providecommand{\shC}{}\renewcommand{\shC}{\mathcal{C}}
\providecommand{\shE}{}\renewcommand{\shE}{\mathcal{E}}
\providecommand{\shJ}{}\renewcommand{\shJ}{\mathcal{J}}
\providecommand{\shK}{}\renewcommand{\shK}{\mathcal{K}}
\providecommand{\shL}{}\renewcommand{\shL}{L}
\providecommand{\shW}{}\renewcommand{\shW}{\mathcal{W}}
\providecommand{\sll}{}\renewcommand{\sll}{\mathfrak{sl}}
\providecommand{\so}{}\renewcommand{\so}{\mathfrak{so}}
\providecommand{\su}{}\renewcommand{\su}{\mathfrak{su}}
\providecommand{\supp}{}\renewcommand{\supp}{\operatorname{supp}}
\providecommand{\tr}{}\renewcommand{\tr}{\operatorname{tr}}
\providecommand{\uu}{}\renewcommand{\uu}{\mathfrak{u}}
\providecommand{\br}{}\renewcommand{\br}[1]{\langle #1\rangle}
\providecommand{\bra}{}\renewcommand{\bra}[1]{\langle #1|}
\providecommand{\file}{}\renewcommand{\file}[1]{\texttt{#1}}
\providecommand{\hp}{}\renewcommand{\hp}[2]{( #1 \,|\, #2 )}
\providecommand{\ip}{}\renewcommand{\ip}[2]{\langle #1, #2\rangle}
\providecommand{\ket}{}\renewcommand{\ket}[1]{|#1\rangle}
\providecommand{\rel}{}\renewcommand{\rel}[2]{D\bigl(#1 \,\|\, #2\bigr)}
\providecommand{\rev}{}\renewcommand{\rev}[1]{\widetilde{#1}}
\providecommand{\slashed}{}\renewcommand{\slashed}[1]{{#1}\mkern-10mu/}
\title{\bfseries The Descent, Lifted: What the pinching discards, a pairing stores: Companion note to T1 and T5: the reversible lift of the descent costs one rung per mode and needs exactly the generator the singularity paper excludes}
\author{Leon Butler\\ \normalsize\href{mailto:lbutler.math@gmail.com}{\texttt{lbutler.math@gmail.com}}\\ \normalsize ORCID: \href{https://orcid.org/0009-0006-4449-5617}{0009-0006-4449-5617}}
\date{Generated by the Transport Lab, 27 October 2026}
\begin{document}
\maketitle
\begin{abstract}
\normalsize
The series' descent, the pinching onto the number-diagonal algebra, is one-way: the root paper proves no homomorphism returns, and the singularity paper shows it sends an entangled pair to a classical shared bit. This note treats the descent as a channel and asks whether it lifts as the transport principle lifts an element without an inverse: reversibly, by adjoining what it discards. It does, at an exact price. On the even (physical) states of $\nu$ modes the descent has Kraus rank $2^{\nu - 1}$, and an even isometric lift needs that many orthogonal even record states; $k$ rungs of the tower supply $2^{k - 1}$, so the minimal record is $\nu$ rungs, one per mode. One lift is explicit: a coupling in the singularity paper's supplied class applied to a record prepared in the Thouless state with all pairing amplitudes one, by a rotation of the family whose member at amplitude $4/3$ is that paper's necessity witness. The pairing may sit on the record or on the pair; no number-conserving dynamics replaces it. The discarded coherence reappears as entanglement across the pair-record cut with the pair's Schmidt coefficients; the classical shared bit is the reduced view of a pure four-mode state. Seen from the two sides of the descent, the one-way door and the exclusion of pairing terms are one fact within the supplied dynamics. The lift does not touch the Hawking angle, and extending the principle from elements to maps is a reading, not a theorem.
\end{abstract}
\medskip
\noindent\textbf{MSC 2020.} Primary 81P40; Secondary 15A66, 46L60, 81R25.\\
\textbf{Keywords.} Conditional expectation; pinching; Stinespring dilation; fermionic parity superselection; Bogoliubov transformation; entanglement; transport principle.
\newpage
\noindent \textbf{Authorship and disclosure. }The founding idea of the series is the author's: that division by zero in a Clifford algebra should be read not as a prohibition but as a transport into a higher-dimensional algebra, and that the transported division need not be accompanied by an inverse operation. The direction of investigation and the critical review of the results are likewise the author's. The mathematical development and the drafting of the text were carried out in extended collaboration with Claude Fable 5 (Anthropic), an AI system.

\noindent\fbox{\parbox{\dimexpr\linewidth-2\fboxsep-2\fboxrule}{\textbf{This is not a deposited paper.} It was generated by the Transport Lab on 27 October 2026 and assembled from approved parts of its store. It follows results across several papers of the series; each is stated with its proof and credited to the paper that first proved it. The official versions of those papers are on Zenodo, in the Transport Papers community, \url{https://zenodo.org/communities/transport-papers}.}}
\newpage

\section{Setting, and what is imported}

\phantomsection\label{c:note.descent-lift.setting-setting-what-imported}
The transport principle is a rule for elements: where an element has no inverse, adjoin the minimum that gives it one, and keep the descent that forgets the adjunction \cite[Thm.~4.13, Prop.~7.5, Thm.~7.6]{ButlerT1}. The descent is a conditional expectation, the pinching onto the algebra that commutes with the ladder's occupation projector, and it is one-way: no homomorphism returns from the ascent to the base. The singularity paper reads the same pinching, mode by mode, as the parity average $\EN(x) = \tfrac12(x + \G x\G)$ and, on several modes, as the pinching onto the joint-number-diagonal algebra \cite[Thm.~5.7, Prop.~9.6]{ButlerT5}; applied to an entangled pair $c\ket{00} + s\ket{11}$ it returns the classical shared bit $\operatorname{diag}(c^{2}, s^{2})$.

This note asks the inverted question. A map without an inverse can still be lifted: enlarged to a reversible map on a bigger system, followed by forgetting the enlargement. That is Stinespring's dilation, and it is the exact analogue for maps of what the transport principle does for elements. The question is whether the descent's lift can be built inside the tower, from rungs, by operations the framework's own premises allow, and what it costs. The answer is computed exactly, in the singularity paper's conventions, and the kill conditions of the section ``The claim, and the kill conditions fixed before computation'' were fixed before the computation.

Three things are fixed on this page. First, the states. The singularity paper's premise (P) says physical states of the modes are even, and its imported univalence rule says the observable algebra is the even subalgebra \cite[Prem.~7.8, Fact~7.9]{ButlerT5}. Everything below is stated on even states, and a single mode then has nothing to lift: its only even pure state is the vacuum, on which the descent is the identity. The descent's first physical action is on the pair. Second, the reduced state. For fermions the reduced state of a subsystem is the restriction of the joint state to the subsystem's algebra; the tensor-product partial trace agrees with it when the complementary factor is even, which is the only case the lift uses, and the script computes the restriction throughout. Third, the vocabulary, for a reader of quantum information who has not read the series: a \emph{rung} is one $(1,1)$ ascent of the tower, which adjoins one fermionic mode \cite[Thm.~4.2]{ButlerT3}; the \emph{descent} is the pinching above; the \emph{supplied class} is the set of generators the singularity paper's completion provides, polynomials in the parity sections $\sigma_{k} = 2N_{k} - 1$, all joint-number-diagonal \cite[Prem.~7.30, Thm.~C.8]{ButlerT5}; a \emph{pairing} generator is $a^{\dagger}_{i}a^{\dagger}_{j}$ or its adjoint, the part of $\mathfrak{so}(2\nu)$ that moves the vacuum \cite[Prop.~13.4]{ButlerT5}; and the two-ledger words \emph{transport-essential}, \emph{convenient} and \emph{absent} grade a result by whether the $(1,1)$ increment does irreplaceable work in it, some work, or none \cite{ButlerT17}.

\begin{assembled}[Imported: the pinching and the one-way door]\label{c:descent-lift.imported-pinching-one-way-door}
For the ascent $B = \Cl(p+1, q+1) \cong \mathrm{Mat}_{2}(A)$ of $A = \Cl(p,q)$ with ladder pair $f, f^{*}$ and occupation projector $P = f^{*}f$, the pinching $\EN(x) = PxP + (1-P)x(1-P)$ is an idempotent conditional expectation onto the commutant of $P$; its kernel is the ladder sector. There is no unital algebra homomorphism $B \to A$: every descent is non-multiplicative, and the pinching is the canonical linear surrogate. \emph{Provenance:} \cite[Prop.~7.5, Thm.~7.6]{ButlerT1}.
\end{assembled}
\begin{proof}
For the ascent $B = \Cl(p+1, q+1) \cong \mathrm{Mat}_{2}(A)$ of $A = \Cl(p,q)$ with ladder pair $f, f^{*}$ and occupation projector $P = f^{*}f$, the pinching $\EN(x) = PxP + (1-P)x(1-P)$ is an idempotent conditional expectation onto the commutant of $P$; its kernel is the ladder sector. There is no unital algebra homomorphism $B \to A$: every descent is non-multiplicative, and the pinching is the canonical linear surrogate. \emph{Provenance:} \cite[Prop.~7.5, Thm.~7.6]{ButlerT1} (\ref{c:transport.pinching-kernel}, \ref{c:transport.door-one-way}).

\end{proof}

\begin{fact}[Imported: the descent on modes, univalence, and the supplied class (external ingredients)]\label{c:descent-lift.imported-descent-modes-univalence-supplied-class-external}
The results from outside the series on which \ref{c:descent-lift.imported-descent-modes-univalence-supplied-class} draws, as cited there: G. C. Wick, A. S. Wightman, E. P. Wigner, The intrinsic parity of elementary particles, Phys. Rev. 88 (1952) 101--105..
\emph{Source:} G. C. Wick, A. S. Wightman, E. P. Wigner, The intrinsic parity of elementary particles, Phys. Rev. 88 (1952) 101--105..
\end{fact}

\begin{assembled}[Imported: the descent on modes, univalence, and the supplied class]\label{c:descent-lift.imported-descent-modes-univalence-supplied-class}
On the mode algebra the descent is the parity average $\EN(x) = \tfrac12(x + \sigma x\sigma)$, and on $\nu$ modes the composite of the mode-wise averages is the pinching onto the joint-number-diagonal algebra. Physical states are even (P), which is the univalence superselection rule: coherent superpositions of opposite fermion parity are not observable, and the observable algebra is the even subalgebra. The completion attaches one ladder pair to each point of the exceptional set and supplies, as generators of the boundary-mode dynamics, the parity sections and their products, all joint-number-diagonal; the exclusion of pairing terms is premised. On the Fock space of $\nu \geq 2$ modes the hops and pair annihilations stabilise the vacuum ray and the pair creations span the tangent of its orbit; the rotation $\ket{00} \mapsto \tfrac35\ket{00} + \tfrac45\ket{11}$ is the pairing rotation on one pair that the envelope corollary needs. The thermofield pair $c\ket{00} + s\ket{11}$ with $\tan\theta = s/c = e^{-\omega/2T}$ has both marginals equal to the Gibbs state $\operatorname{diag}(c^{2}, s^{2})$; the sample point $e^{-\omega/T} = \tfrac{9}{16}$ gives $(c, s) = (\tfrac45, \tfrac35)$. \emph{Provenance:} \cite[Thm.~5.7, Prop.~9.6, Prem.~7.8, Fact~7.9, Prem.~7.30, Thm.~C.8, Prop.~13.4, Cor.~11.6, Lem.~7.2, Thm.~7.27]{ButlerT5}; the univalence rule \cite{WWW1952}.
\end{assembled}
\begin{proof}
On the mode algebra the descent is the parity average $\EN(x) = \tfrac12(x + \sigma x\sigma)$, and on $\nu$ modes the composite of the mode-wise averages is the pinching onto the joint-number-diagonal algebra. Physical states are even (P, \ref{c:singularity.p-univalence}), which is the univalence superselection rule: coherent superpositions of opposite fermion parity are not observable, and the observable algebra is the even subalgebra. The completion attaches one ladder pair to each point of the exceptional set and supplies, as generators of the boundary-mode dynamics, the parity sections and their products, all joint-number-diagonal; the exclusion of pairing terms is premised (\ref{c:singularity.l-boundary-mode-locality}). On the Fock space of $\nu \geq 2$ modes the hops and pair annihilations stabilise the vacuum ray and the pair creations span the tangent of its orbit; the rotation $\ket{00} \mapsto \tfrac35\ket{00} + \tfrac45\ket{11}$ is the pairing rotation on one pair that the envelope corollary needs. The thermofield pair $c\ket{00} + s\ket{11}$ with $\tan\theta = s/c = e^{-\omega/2T}$ has both marginals equal to the Gibbs state $\operatorname{diag}(c^{2}, s^{2})$; the sample point $e^{-\omega/T} = \tfrac{9}{16}$ gives $(c, s) = (\tfrac45, \tfrac35)$. \emph{Provenance:} \cite[Thm.~5.7, Prop.~9.6, Prem.~7.8, Fact~7.9, Prem.~7.30, Thm.~C.8, Prop.~13.4, Cor.~11.6, Lem.~7.2, Thm.~7.27]{ButlerT5} (\ref{c:singularity.envelope-endpoint-why-boundary-mode-locality}, \ref{c:singularity.imported-univalence-superselection-rule}, \ref{c:singularity.p-reduced-metric-level-reading-many}, \ref{c:singularity.l-stabiliser-tangent}, \ref{c:singularity.descent-parity-average}, \ref{c:singularity.purity-output-partner-commutant}, \ref{c:singularity.schmidt-forcing-proved-outright}); the univalence rule \cite{WWW1952}.

This restates \ref{c:descent-lift.imported-descent-modes-univalence-supplied-class-external}, \ref{c:singularity.class-relative-no-go-no-supplied}.

\end{proof}

\begin{fact}[Imported: Stinespring, Kraus, Choi]\label{c:descent-lift.imported-stinespring-kraus-choi}
A completely positive trace-preserving map $\Phi$ on the operators of a finite-dimensional space $H$ has a Kraus form $\Phi(\rho) = \sum_{k} K_{k}\rho K_{k}^{\dagger}$, and a dilation: an isometry $W : H \to H \otimes H_{R}$ with $\Phi(\rho) = \Tr_{R}(W\rho W^{\dagger})$, given by $W\psi = \sum_{k} K_{k}\psi \otimes \ket{k}$ for an orthonormal family $\ket{k}$ of $H_{R}$. The minimal number of Kraus operators, and the minimal dimension of $H_{R}$, is the rank of the Choi matrix $\sum_{ij}\Phi(\ket{i}\bra{j}) \otimes \ket{i}\bra{j}$; every dilation is the minimal one followed by an isometry on the environment, so every dilation has the form $W = \sum_{k} K_{k} \otimes \ket{r_{k}}$ with the $\ket{r_{k}}$ orthonormal and the $K_{k}$ a Kraus family of $\Phi$. \emph{Provenance:} \cite{Stinespring1955, Kraus1983, Choi1975}.
\emph{Source:} W. F. Stinespring, Positive functions on $C^*$-algebras, Proc. Amer. Math. Soc. 6 (1955) 211--216.; K. Kraus, States, Effects, and Operations: Fundamental Notions of Quantum Theory, Lecture Notes in Physics 190, Springer (1983).; M.-D. Choi, Completely positive linear maps on complex matrices, Linear Algebra Appl. 10 (1975) 285--290..
\end{fact}

\begin{fact}[Imported: fermionic operations are even, and reduced states are restrictions]\label{c:descent-lift.imported-fermionic-operations-are-even-reduced}
In a fermionic system with the univalence rule, the physically implementable unitaries are those commuting with the total parity, the even ones; an odd operation would carry an even state to an odd one or to a superposition of the two. The reduced state of a subsystem is the restriction of the state to the subsystem's algebra; when the complementary subsystem's states are even, this agrees with the partial trace in the tensor-product basis, and a product state $\omega_{S} \otimes \omega_{R}$ exists whenever one factor is even. \emph{Provenance:} \cite{BravyiKitaev2002}; \cite[Thm.~1]{ArakiMoriya2003}; the univalence rule \cite{WWW1952}, imported by the singularity paper as Fact \ref{c:descent-lift.imported-descent-modes-univalence-supplied-class}.
\emph{Source:} S. B. Bravyi, A. Yu. Kitaev, Fermionic quantum computation, Ann. Phys. 298 (2002) 210--226.; H. Araki, H. Moriya, Joint extension of states of subsystems for a CAR system, Commun. Math. Phys. 237 (2003) 105--122., Thm.~1; G. C. Wick, A. S. Wightman, E. P. Wigner, The intrinsic parity of elementary particles, Phys. Rev. 88 (1952) 101--105..
\end{fact}

\begin{fact}[Imported: Thouless states and Witt vacua (external ingredients)]\label{c:descent-lift.imported-thouless-states-witt-vacua-external}
The results from outside the series on which \ref{c:descent-lift.imported-thouless-states-witt-vacua} draws, as cited there: D. J. Thouless, Stability conditions and nuclear rotations in the Hartree-Fock theory, Nucl. Phys. 21 (1960) 225--232..
\emph{Source:} D. J. Thouless, Stability conditions and nuclear rotations in the Hartree-Fock theory, Nucl. Phys. 21 (1960) 225--232..
\end{fact}

\begin{assembled}[Imported: Thouless states and Witt vacua]\label{c:descent-lift.imported-thouless-states-witt-vacua}
For an antisymmetric matrix $A$ the state $\exp(\sum_{i<j}A_{ij}a^{\dagger}_{i}a^{\dagger}_{j})\ket{0}$ is the vacuum of the Witt decomposition obtained from the standard one by the corresponding Bogoliubov rotation; on two modes it is $\ket{00} + A_{12}\ket{11}$, and every even state vector of at most three modes is a Witt vacuum (pure, by Cartan's theorem), of this exponential form exactly when its overlap with $\ket{0}$ is nonzero ($\ket{11}$ on two modes is a Witt vacuum not of this form). \emph{Provenance:} \cite{Thouless1960}; Cartan's theorem as imported by the laboratory paper \cite[Fact~2.10]{ButlerT16}.
\end{assembled}
\begin{proof}
For an antisymmetric matrix $A$ the state $\exp(\sum_{i<j}A_{ij}a^{\dagger}_{i}a^{\dagger}_{j})\ket{0}$ is the vacuum of the Witt decomposition obtained from the standard one by the corresponding Bogoliubov rotation; on two modes it is $\ket{00} + A_{12}\ket{11}$, and every even state vector of at most three modes is a Witt vacuum (pure, by Cartan's theorem), of this exponential form exactly when its overlap with $\ket{0}$ is nonzero ($\ket{11}$ on two modes is a Witt vacuum not of this form). \emph{Provenance:} \cite{Thouless1960}; Cartan's theorem as imported by the laboratory paper \cite[Fact~2.10]{ButlerT16} (\ref{c:one-relation.imported-cartans-purity-theorem-thouless-form}).

This restates \ref{c:descent-lift.imported-thouless-states-witt-vacua-external}.

\end{proof}

\section*{What we assume}
The results below are taken as given in this paper; each result that rests on one says so.
\begin{postulate}[(L$_E$): boundary-mode locality, as descent covariance on finite windows]\label{c:singularity.l-boundary-mode-locality}
\emph{Premise (L$_{E}$), boundary-mode locality.} On every finite window of the boundary modes, $v_{1}, \dots, v_{n}$ satisfying the CAR, the internal dynamics $\alpha_{t}$ of the boundary modes commutes with the descent at all times:
\[
\alpha_{t}\circ\mathcal{E} \;=\; \mathcal{E}\circ\alpha_{t} \quad \text{for all } t, \qquad \mathcal{E} = E_{1}\circ\cdots\circ E_{n}, \quad E_{k}(x) = N_{k}xN_{k} + (1 - N_{k})x(1 - N_{k}),
\]
$\mathcal{E}$ the descent of the window (Proposition \emph{The descent, and what it erases} [T5, Prop 5.5]).

\emph{The version 1 wording, equivalent.} By Theorem \emph{On finite windows, boundary-mode locality is equivalent to the dynamics commuting with the descent for all times}, on every finite window this holds if and only if the dynamics is generated by elements of the abelian algebra $\mathcal{D}$ that the $\sigma_{k} = 2N_{k} - 1$ generate, diagonal in the joint number basis: the internal dynamics is generated by the geometric data the framework supplies on $E$, the boost bivector field $x \mapsto \sigma_{r}(x)$ and its products. Cross-point products such as $\sigma_{j}\sigma_{k}$ correlate phases without moving occupation; hopping $v_{j}^{*}v_{k} + v_{k}^{*}v_{j}$, pairing $v_{j}^{*}v_{k}^{*} + v_{k}v_{j}$ and every other generator with a component outside $\mathcal{D}$ are excluded by the one clause. Within the class of generators the framework supplies, the exclusion of hopping is a theorem (Theorem \ref{c:singularity.class-relative-no-go-no-supplied}); what the premise adds is the exclusion of unsupplied generators, pairing terms and the possible differential dynamics, which Proposition \ref{c:singularity.l-stabiliser-tangent} places in the tangent of the vacuum orbit and in its stabiliser.

\emph{What it forbids.} By Theorem \emph{Number-diagonal dynamics carries operator support only to directly coupled modes, with no exponential growth, so no dynamics (L\_E) admits realises the kappa\^{}{-1} ln N spreading law in real time}, every dynamics the premise admits carries an evolved mode operator only to the modes directly coupled to it, and every local commutator and out-of-time-order correlator is almost periodic and bounded: the premise forbids exponential real-time commutator growth. The real-time spreading law $t_{*} = \kappa^{-1}\ln N$ of Proposition \emph{The spreading law, and the chaos bound met, not saturated} [T8, Asm 6.1] is therefore not realised under it, and the scrambling time is a modular-time statement: the growth $e^{\kappa t}$ is a rate of the rotor flow, the KMS-strip continuation of the phase flow that is the windows' time evolution (Theorem \emph{The phase flow is the time evolution; the rotor flow is its KMS continuation}). The premise is kept over fast real-time scrambling; the two cannot both hold.
\end{postulate}

\begin{postulate}[(P): univalence]\label{c:singularity.p-univalence}
Physical states of the boundary modes are even: $\omega \circ \alpha = \omega$, equivalently $\omega$ annihilates the odd sector $\Aod$.
\end{postulate}

\section{The claim, and the kill conditions fixed before computation}

\phantomsection\label{c:note.descent-lift.setting-claim-kill-conditions-fixed-before}
The kill conditions and predictions below are as registered, in the deposited pre-registration \cite{ButlerPreregT05dilation}; every departure between that wording and what the note does is collected in Remark \ref{c:note.descent-lift.departures-registered-wording}.

\phantomsection\label{c:note.descent-lift.setting-claim-kill-conditions-fixed-before-2}
\textbf{Claim D (the descent lifts, at a price).} The descent has a lift within the tower: an even isometric dilation whose record is a set of rungs, realised by an even unitary that factorises as a coupling in the singularity paper's supplied class composed with a pairing rotation on the record. The coherence the descent discards is stored as system-record entanglement with Schmidt coefficients equal to the system's amplitudes. The only step outside the supplied class is the pairing rotation, so within the dynamics the framework supplies the descent has no lift: the root paper's one-way door and the singularity paper's exclusion of pairing terms are one fact.

\begin{itemize}
\item[K1] (one rung suffices) If an even isometric dilation of the single-mode parity average exists with a one-rung record, the count ``two rungs for one mode'' is dead. Test: the Kraus operators span $\{1, \G\}$, both even and number-diagonal, so an even isometry has the form $P_{0} \otimes r_{0} + P_{1} \otimes r_{1}$ with $r_{0}, r_{1}$ even record states that must be orthogonal ($\EN$ kills the coherence); the even sector of one rung is one-dimensional.
\item[K2] (no pairing needed) If an even unitary generated by number-conserving elements only (hops within the record, hops between system and record, Cartan elements, joint-number-diagonal couplings) realises a lift from the record vacuum, the clause ``the only step outside the supplied class is a pairing rotation'' is dead. Test: such a unitary conserves the total number, the joint state stays in total number $\leq \nu$, and the record's even states there are too few to be orthogonal.
\item[K3] (coupling outside the supplied class) If the controlled coupling $C$ of the factorisation is not joint-number-diagonal (not a polynomial in the parity sections), the factorisation clause is dead.
\item[K4] (coherence not stored) If the lifted state of $c\ket{0} + s\ket{1}$ does not have reduced spectrum $(c^{2}, s^{2})$ on the system and on the record, with the joint state pure, the reading ``the discarded coherence is the system-record entanglement'' is dead.
\item[K5] (the principle) The root paper's transport principle concerns elements without inverses. Applying it to a map without an inverse is an extension. If nothing in the root paper forces a non-invertible map to be lifted, the lift is a reading and the result is transport-convenient at most. Expected to trigger; recorded as such, not as a kill of the mathematics.
\end{itemize}
Predictions, as registered: the minimal even record for one mode is two rungs, in their even sector, and for $\nu$ modes $\nu + 1$ rungs, the per-mode product lift using $2\nu$ and not being minimal (P1); the record must leave its vacuum by two quanta, a pairing creation, and the stabiliser of the record vacuum fixes its ray (P2); $C = 1 - 2P^{R}_{11}P^{S}_{1}$ for one mode (P3); at $(c, s) = (\tfrac45, \tfrac35)$ the lifted state is $c\ket{0}\ket{+} + s\ket{1}\ket{-}$ with reduced spectra $(\tfrac{16}{25}, \tfrac{9}{25})$ on both sides, the system marginal being the thermal marginal at $e^{-\omega/T} = \tfrac{9}{16}$ (P4); the record preparation is the necessity witness rotation at amplitude $1$, and with amplitude $A$ the lift is a partial descent with coherence factor $(1 - A^{2})/(1 + A^{2})$ (P5); a one-rung stochastic version exists and records nothing recoverable (P6); K5 triggers, with ledger entry transport-convenient (P7).

\section{The descent as a channel on physical states}

\begin{proposition}[Kraus form of the descent on even states]\label{c:descent-lift.kraus-form-descent-even-states}
Let $\G_{k} = (-1)^{N_{k}}$ be the parity of mode $k$ and $\G^{\kappa} = \prod_{k}\G_{k}^{\kappa_{k}}$ the parity monomials, $\kappa \in \{0,1\}^{\nu}$. On the full Fock space of $\nu$ modes the descent is $\EN(\rho) = 2^{-\nu}\sum_{\kappa}\G^{\kappa}\rho\,\G^{\kappa}$, of Kraus rank $2^{\nu}$. On the even sector, where $\G_{\mathrm{tot}} = \prod_{k}\G_{k} = 1$, the monomials pair up and the descent is $\EN(\rho) = 2^{1-\nu}\sum_{\kappa_{\nu} = 0}\G^{\kappa}\rho\,\G^{\kappa}$, of Kraus rank $2^{\nu - 1}$, every Kraus operator even and number-diagonal. For $\nu = 1$ it is the identity. For $\nu = 2$ it is $\EN(\rho) = \tfrac12(\rho + \G_{0}\rho\,\G_{0})$ on $\operatorname{span}\{\ket{00}, \ket{11}\}$, which kills exactly the pairing coherence: $c\ket{00} + s\ket{11} \mapsto \operatorname{diag}(c^{2}, s^{2})$.
\end{proposition}
\begin{proof}
The mode-wise parity average is $\tfrac12(\rho + \G_{k}\rho\G_{k})$ and the averages commute, so their composite is the stated sum (Fact \ref{c:descent-lift.imported-descent-modes-univalence-supplied-class}); it kills every matrix unit $\ket{A}\bra{B}$ with $A \neq B$, since some $\G_{k}$ has opposite signs on $A$ and $B$, and fixes the diagonal ones: the pinching. The Kraus rank is the Choi rank (\ref{c:descent-lift.imported-stinespring-kraus-choi}), which the script computes as $4$ on the full two-mode space and $2$ on its even sector, and $4$ on the even sector of three modes. On the even sector $\G_{\nu} = \prod_{k < \nu}\G_{k}$, so the $2^{\nu}$ monomials collapse to $2^{\nu-1}$ distinct operators, each appearing twice. On one mode the even sector is the vacuum line. On the pair, $\G_{0}$ and $\G_{1}$ agree on the even sector ($+1$ on $\ket{00}$, $-1$ on $\ket{11}$), and the script checks the descent of $(\tfrac45)\ket{00} + (\tfrac35)\ket{11}$ is $\operatorname{diag}(\tfrac{16}{25}, \tfrac{9}{25})$ with the coherence $\tfrac{12}{25}$ removed.

\end{proof}

\begin{remark}[Where the descent first acts]\label{c:note.descent-lift.where-descent-first-acts}
On a single mode the descent acts only on the odd-even coherence $\ket{0}\bra{1}$, which (P) declares unphysical: the single-mode descent is the identity on every physical state. The descent's first physical action is on the pair, and what it removes there is the pairing coherence between $\ket{00}$ and $\ket{11}$, the coherence that makes the pair entangled. That is the entangled-pair-to-classical-bit computation of the singularity paper's descent, and it is the object lifted below.
\end{remark}

\section{The lift, and its cost}

\begin{definition}[Lift]\label{c:obj.descent-lift.lift}
A \emph{lift} of the descent on $\nu$ modes is an isometry $W : H_{S} \to H_{S} \otimes H_{R}$, $H_{R}$ the Fock space of $k$ rungs, that is even, $W\G_{S} = (\G_{S}\otimes\G_{R})W$, and induces the descent on the even sector: the restriction of $W\rho W^{\dagger}$ to the system's algebra equals $\EN(\rho)$ for every even $\rho$. A lift is \emph{realised} by an even unitary $U$ on $H_{S}\otimes H_{R}$ with $W = U(\,\cdot \otimes \ket{0_{R}})$.
\end{definition}

\begin{theorem}[The lift costs one rung per mode, and needs a pairing]\label{c:descent-lift.lift-costs-one-rung-per-mode}
Let $\nu \geq 2$.
\begin{enumerate}
\item[(i)] A lift has the form $W\ket{A} = \ket{A} \otimes r_{A}$ on the even basis states $\ket{A}$, with $r_{A} = \sum_{\kappa}(K_{\kappa})_{AA}\ket{r_{\kappa}}$ for orthonormal record states $\ket{r_{\kappa}}$, one per Kraus operator, and every $\ket{r_{\kappa}}$ even. Hence the record must contain $2^{\nu-1}$ orthonormal even states.
\item[(ii)] The even sector of $k$ rungs has dimension $2^{k-1}$. Hence a lift needs $k \geq \nu$ rungs, one per mode, and one rung never suffices.
\item[(iii)] With $\nu$ rungs a lift exists: $W = \sum_{\kappa_{\nu} = 0}\G^{\kappa} \otimes \ket{r_{\kappa}}$, the $\ket{r_{\kappa}}$ any orthonormal basis of the record's even sector; for the pair, $W = 1 \otimes \ket{00} + \G_{0} \otimes \ket{11}$, up to the normalisation $W^{\dagger}W = 2$.
\item[(iv)] The lift factorises as $W = C\,(\,\cdot \otimes \varphi)$ with $\varphi = \sum_{\kappa}\ket{r_{\kappa}}$ and $C = \sum_{\kappa}\G^{\kappa} \otimes \ket{r_{\kappa}}\bra{r_{\kappa}}$. For the pair, $C = 1 - 2P^{R}_{11}N_{0}$, a polynomial in the parity sections, joint-number-diagonal, even, an involution: an element of the supplied class. And $\varphi = \ket{00} + \ket{11}$ is the Thouless state of the record pair with amplitude $1$, a Witt vacuum of a rotated decomposition.
\item[(v)] Any even unitary $U$ realising a lift from the record vacuum has a generator that does not conserve the total number; among the quadratic generators, those of the Bogoliubov algebra $\mathfrak{so}(2(\nu + k))$, that is a generator with a nonzero pairing component. A unitary generated by number-conserving elements (hops within the record, hops between system and record, Cartan elements, joint-number-diagonal couplings) realises no lift. The pairing may sit on the record, as in (iv), or on the system and be carried to the record by hops.
\end{enumerate}
\end{theorem}
\begin{proof}
(i) By Fact \ref{c:descent-lift.imported-stinespring-kraus-choi} every dilation is $W = \sum_{\kappa}K_{\kappa} \otimes \ket{r_{\kappa}}$ with $K_{\kappa}$ a Kraus family and the $\ket{r_{\kappa}}$ orthonormal, and by Proposition \ref{c:descent-lift.kraus-form-descent-even-states} the Kraus operators are diagonal in the number basis, so $W\ket{A} = \ket{A} \otimes r_{A}$. Evenness of $W$ and of $\ket{A}$ force $r_{A}$ even for every $A$; since the matrix $((K_{\kappa})_{AA})$ is invertible (the Kraus operators are linearly independent), each $\ket{r_{\kappa}}$ is a combination of the $r_{A}$ and is even. The record must therefore contain $2^{\nu-1}$ orthonormal even states. (ii) The even basis states of $k$ rungs are the subsets of even size, $2^{k-1}$ of them; the script counts them for $k \leq 4$. (iii) $W$ is even because each $\G^{\kappa}$ and each $\ket{r_{\kappa}}$ is; $W^{\dagger}W = \sum_{\kappa}(\G^{\kappa})^{2} = 2^{\nu-1}$; and the restriction of $W\rho W^{\dagger}$ to the system's algebra is $\sum_{\kappa}\G^{\kappa}\rho\G^{\kappa}$ because the record states are orthonormal and even, which is $2^{\nu-1}\EN(\rho)$ on even $\rho$. The script verifies evenness, the norm and the induced channel on every even matrix unit for $\nu = 2$ (two rungs) and $\nu = 3$ (three rungs), computing the restriction to the system's algebra directly and, for the pair, checking that it agrees with the tensor-product partial trace. (iv) $C(\ket{A} \otimes \varphi) = \sum_{\kappa}\G^{\kappa}\ket{A} \otimes \ket{r_{\kappa}} = W\ket{A}$. For the pair, $\G_{0} = 1 - 2N_{0}$ gives $C = 1 - 2P^{R}_{11}N_{0} = 1 - \tfrac14(1 - \G_{0})(1 - \G_{2})(1 - \G_{3})$; the script checks the polynomial identity, diagonality, commutation with every $N_{k}$ and with the total parity, and $C^{2} = 1$. That $\varphi$ is a Witt vacuum is Fact \ref{c:descent-lift.imported-thouless-states-witt-vacua}, checked by the $M^{2} = -N^{4}I$ criterion on the correlation matrix. (v) Let $U$ be generated by number-conserving elements; then $U$ commutes with the total number. By (i) a lift keeps the system in $\ket{A}$ while writing $r_{A}$, so $U(\ket{A} \otimes \ket{0_{R}})$ has total number $|A|$ and system part $\ket{A}$, whence the record part has number $0$: $r_{A} = \ket{0_{R}}$ for every $A$, the $r_{A}$ are not orthogonal, and the induced map keeps the coherence. The script lists ten number-conserving generators, checks each commutes with the total number, and enumerates the joint basis states of the required number and parity: for the pair, both $A = \ket{00}$ and $A = \ket{11}$ admit only the record state $\ket{00}$. A record that reaches $\ket{11}$ from $\ket{00}$ while the system's number is unchanged has changed its own number by two, which, within the record's own quadratic algebra, is the pair creation $a^{\dagger}_{2}a^{\dagger}_{3}$; the script checks that every stabiliser element of the record vacuum (hops, the annihilation, the Cartan elements) annihilates $\ket{00}$ (Fact \ref{c:descent-lift.imported-descent-modes-univalence-supplied-class}), and that the rational pairing rotor $(1 - P) + \tfrac35 P + \tfrac45 X$, with $X = a^{\dagger}_{2}a^{\dagger}_{3} - a_{3}a_{2}$ and $P = -X^{2}$ the projection onto the record's even sector $\operatorname{span}\{\ket{00}, \ket{11}\}$ (the element $\tfrac35 + \tfrac45 X$ alone is orthogonal only on that sector), is unitary, even, and carries $\ket{00}$ to $\tfrac35\ket{00} + \tfrac45\ket{11}$, the singularity paper's necessity witness; $\varphi$ is the member of the same family at $\cos\theta = \sin\theta$. The pairing need not act on the record. The script checks that the system's pairing rotor at $\cos\theta = \sin\theta$, unnormalised, $1 + (a^{\dagger}_{0}a^{\dagger}_{1} - a_{1}a_{0})$, conjugated by the two system-record swaps (beam splitters at $(\cos, \sin) = (0, 1)$ between modes $0, 2$ and $1, 3$) and followed by $C$, is an even map inducing $2\EN$ on the even sector: a lift whose only non-number-conserving generator acts on the system. Hops move the pairing between system and record; nothing removes it.

\end{proof}

\begin{corollary}[Parity alone forces the second rung]\label{c:descent-lift.parity-alone-forces-second-rung}
For the pair, both one-rung candidates exist as maps: Stinespring's minimal dilation $\psi \otimes \ket{0} + \G_{0}\psi \otimes \ket{1}$ and the occupation copy $P_{00}\psi \otimes \ket{0} + P_{11}\psi \otimes \ket{1}$ each induce the descent on the even sector. Both are odd: the first because its record states $\ket{0} \pm \ket{1}$ mix the parities, the second because it sends the even state $\ket{11}$ to a joint state of odd parity. Neither is an operation the univalence rule allows, and there is no other obstruction.
\end{corollary}
\begin{proof}
The script checks the induced map on the even sector's four matrix units for both, and the failure of evenness for both, by evaluating $W\G_{S}$ against $(\G_{S}\otimes\G_{R})W$ on the even basis. That no even one-rung lift exists is Theorem \ref{c:descent-lift.lift-costs-one-rung-per-mode}(i) and (ii): a lift needs two orthonormal even record states and one rung has one. That parity is the only obstruction is shown by the two-rung occupation copy $P_{00}\psi \otimes \ket{00} + P_{11}\psi \otimes \ket{11}$, which is even and is the lift of Theorem \ref{c:descent-lift.lift-costs-one-rung-per-mode}(iii) written in the record basis $\ket{00} \pm \ket{11}$: $W = P_{00} \otimes (\ket{00} + \ket{11}) + P_{11} \otimes (\ket{00} - \ket{11})$; the script checks the identity, the evenness and the induced channel.

\end{proof}

\begin{remark}[Reading]\label{c:note.descent-lift.reading}
The tower adjoins one mode per rung and one qubit per rung \cite[Thm.~4.2]{ButlerT3}. The theorem says the descent's reversible lift costs exactly that: one rung per mode of the system, and the count is forced by univalence, not by the dimension of the environment, which Stinespring alone would put at $2^{\nu-1}$ states of any kind, $\nu - 1$ qubits, one rung fewer. What the second rung buys, for the pair, is a two-level record that is the even sector of a mode pair rather than a single mode. The record's preparation is the one step the framework's supplied dynamics does not contain, and it is the same rotation the singularity paper had to assume in order to return its boundary modes.
\end{remark}

\section{What the lift stores}

\begin{corollary}[The classical shared bit is the reduced view of a pure four-mode state]\label{c:descent-lift.classical-shared-bit-reduced-view-pure}
Let the pair be in the thermofield state $c\ket{00} + s\ket{11}$, $(c, s) = (\tfrac45, \tfrac35)$. Its lift is the pure, even four-mode state
\[
c\,\ket{00} \otimes (\ket{00} + \ket{11}) + s\,\ket{11} \otimes (\ket{00} - \ket{11}),
\]
up to normalisation. The pair's marginal is $\operatorname{diag}(\tfrac{16}{25}, \tfrac{9}{25})$ on $\{\ket{00}, \ket{11}\}$, exactly the descent of the pair, the classical shared bit; each single mode of the pair keeps the thermal marginal $\operatorname{diag}(\tfrac{16}{25}, \tfrac{9}{25})$ of the singularity paper's sample point; and the record's marginal has the same spectrum $(\tfrac{16}{25}, \tfrac{9}{25})$. The entanglement of the pair, with Schmidt coefficients $(c, s)$ across the mode-mode cut, has become entanglement with the same Schmidt coefficients across the pair-record cut.
\end{corollary}
\begin{proof}
For the two-rung lift $W$ of \ref{c:descent-lift.lift-costs-one-rung-per-mode}(iii), $W(c\ket{00} + s\ket{11}) = c\ket{00}\otimes\ket{00} + c\ket{00}\otimes\ket{11} + s\ket{11}\otimes\ket{00} - s\ket{11}\otimes\ket{11}$ since $\G_{0}\ket{11} = -\ket{11}$; regrouped, it is the display. The script checks the joint state is even, has norm two, restricts to $\operatorname{diag}(\tfrac{16}{25}, \tfrac{9}{25})$ on the pair's algebra and on each mode's algebra, and that the record's marginal, supported on its even sector, has trace $1$ and determinant $\tfrac{144}{625}$, whose roots are $\tfrac{16}{25}$ and $\tfrac{9}{25}$.

\end{proof}

\begin{corollary}[Partial lifts, and the witness overshoots]\label{c:descent-lift.partial-lifts-witness-overshoots}
With the record prepared in $\ket{00} + A\ket{11}$ instead of $\varphi$, the same coupling induces the partial descent $\rho \mapsto (\rho + A^{2}\G_{0}\rho\G_{0})/(1 + A^{2})$, which leaves the fraction $(1 - A^{2})/(1 + A^{2})$ of the pair's coherence: all of it at $A = 0$, $\tfrac{7}{25}$ at $A = \tfrac34$, none at $A = 1$, and $-\tfrac{7}{25}$ at $A = \tfrac43$, the amplitude of the singularity paper's witness rotation, which overshoots the descent into a sign flip. The descent is complete if and only if the record is maximally paired.
\end{corollary}
\begin{proof}
For the coupling $C$ of \ref{c:descent-lift.lift-costs-one-rung-per-mode}(iv), $C(\psi \otimes (\ket{00} + A\ket{11})) = \psi\otimes\ket{00} + A\,\G_{0}\psi\otimes\ket{11}$, and the restriction to the pair's algebra is $\rho + A^{2}\G_{0}\rho\G_{0}$ up to normalisation; the coherence of $\rho$ is multiplied by $(1 - A^{2})/(1 + A^{2})$. The script computes the four values.

\end{proof}

\begin{proposition}[A classical record reproduces the descent and stores nothing]\label{c:descent-lift.classical-record-reproduces-descent-stores-nothing}
With the record pair in the even mixture $\tfrac12(\ket{00}\bra{00} + \ket{11}\bra{11})$ and the same coupling $C$, the induced map on the pair is again the descent, but the joint output of a pure input has purity $\tfrac12$, against $1$ for the lift. One rung suffices for this: a single rung in the maximally mixed state $\tfrac12(\ket{0}\bra{0} + \ket{1}\bra{1})$, with the coupling $1 - 2P^{R}_{1}N_{0}$, also induces the descent, its joint output commuting with the total parity and again of purity $\tfrac12$. Two rungs are needed only for the coherent lift.
\end{proposition}
\begin{proof}
$C$, the coupling of \ref{c:descent-lift.lift-costs-one-rung-per-mode}(iv), acts on the two branches as $1$ and as $\G_{0}$, so the induced map is $\tfrac12(\rho + \G_{0}\rho\G_{0})$; the joint output is a mixture of two orthogonal pure states with equal weights. The one-rung version is the same computation on the eight-dimensional space: the record is a mixture of the two parities, which the univalence rule allows, and its coupling is again a polynomial in the parity sections. The script forms both joint density matrices and computes their restrictions, their purities and, for the one-rung case, the commutation with the total parity.

\end{proof}

\section{What this does to the two papers, and what it does not do}

\begin{proposition}[Within the supplied dynamics the descent has no lift]\label{c:descent-lift.within-supplied-dynamics-descent-has-no}
Let the admissible generators on the system and any number of rungs be the Hermitian elements of the singularity paper's supplied class together with the Hermitian quadratic elements, in $\mathfrak{so}(2(\nu + k))$, that stabilise the joint vacuum ray: all hops and Cartan elements, on the system, on the rungs and between them. Then no unitary generated by admissible elements realises a lift of the descent from the rungs' vacuum. A lift exists in the ascent by $\nu$ rungs when a pairing rotation is added to the admissible generators, and never without a generator that fails to conserve the total number; the lift of Theorem \ref{c:descent-lift.lift-costs-one-rung-per-mode}(iv) places the pairing on the rungs.
\end{proposition}
\begin{proof}
A Hermitian quadratic element is $X = \sum_{i,j} h_{ij}a^{\dagger}_{i}a_{j} + \sum_{i<j} (p_{ij}a^{\dagger}_{i}a^{\dagger}_{j} + \bar{p}_{ij}a_{j}a_{i}) + \lambda$, and $X\ket{0} = \sum_{i<j} p_{ij}a^{\dagger}_{i}a^{\dagger}_{j}\ket{0} + \lambda\ket{0}$ with the states $a^{\dagger}_{i}a^{\dagger}_{j}\ket{0}$ independent; if $X$ stabilises the vacuum ray then $p = 0$, and $X$ is a combination of hops and Cartan elements, which conserve the total number. (The stabiliser of the ray in the complexified algebra also contains the pair annihilations, which are not Hermitian; the script checks that the Hermitian element $a_{2}a_{3} + a^{\dagger}_{3}a^{\dagger}_{2}$ built from one carries the vacuum to $\ket{11}$.) Every admissible generator therefore conserves the total number, since the supplied class is joint-number-diagonal (\ref{c:descent-lift.imported-descent-modes-univalence-supplied-class}), and Theorem \ref{c:descent-lift.lift-costs-one-rung-per-mode}(v) excludes a lift; the same theorem gives the necessity clause. Adding the pairing rotation gives Theorem \ref{c:descent-lift.lift-costs-one-rung-per-mode}(iii) and (iv). The restriction to quadratic elements is needed: the quartic element $I_{4}N_{0}(a^{\dagger}_{2}a^{\dagger}_{3} - a_{3}a_{2})$ is Hermitian, even, and kills the joint vacuum, but does not conserve the total number, and its rotation at $\pi/2$ is the unitary $(1 - N_{0}) + N_{0}(P^{R}_{\mathrm{odd}} + a^{\dagger}_{2}a^{\dagger}_{3} - a_{3}a_{2})$, which carries $\psi \otimes \ket{00}$ to the two-rung occupation copy of Corollary \ref{c:descent-lift.parity-alone-forces-second-rung}, a lift. Here the imaginary unit is $I_{4}$, which commutes with every mode operator of the rungs and with the number operator $N_{0}$, and $I_{4}^{*} = -I_{4}$, so $I_{4}$ times the anti-Hermitian $N_{0}(a^{\dagger}_{2}a^{\dagger}_{3} - a_{3}a_{2})$ is Hermitian. It is a pairing on the rungs conditioned on the system's occupation; the script checks each of these statements, the unitary being entered in closed form. What every lift needs is the non-conservation of number; among quadratic generators that is a pairing.

\end{proof}

\phantomsection\label{c:note.descent-lift.setting-what-does-two-papers-what}
\textbf{One fact.} The root paper proves that the descent is one-way at the level of algebras: no homomorphism returns. The singularity paper premises that the completion supplies no generator off-diagonal in the joint number basis, and its own envelope corollary shows that one pairing rotation is what the return of its boundary modes needs. Proposition \ref{c:descent-lift.within-supplied-dynamics-descent-has-no} joins the two: inside the dynamics the framework supplies, the descent cannot be lifted, and the missing ingredient is a pairing rotation, on the adjoined rungs or on the modes themselves. The one-way door and the pairing exclusion are the same statement seen from the two sides of the descent.

\phantomsection\label{c:note.descent-lift.setting-what-does-two-papers-what-2}
\textbf{Whether the principle produces the pairing.} A second registered experiment, deposited as \cite{ButlerPreregT05pairing}, asks the converse question: does the root paper's own procedure, run on the pairing creator, produce the rotation that T5 excludes? The root paper has no element-level completion of a non-invertible element: its no-go theorem closes that route for null vectors, its type-N repair acts on points, and its type-P rule is that nothing is adjoined \cite[Thm.~4.2, Prop.~6.1]{ButlerT1}. What can be run is the six-step algorithm \cite[Table~1]{ButlerT1}.

\begin{assembled}[Imported: the taxonomy, the type-P rule, and the algorithm]\label{c:descent-lift.imported-taxonomy-type-p-rule-algorithm}
Every zero divisor of a nondegenerate real Clifford algebra is annihilated on each side by a nontrivial idempotent; in a simple algebra the idempotent is non-central and the obstruction is type P, refined to type N exactly when the annihilating idempotent is $wv$ for a hyperbolic pair $(v, w)$ of null vectors, $vw + wv = 1$. Step 2 of the algorithm diagnoses by three tests in order: a vector of square zero is type N; a central idempotent is type C; otherwise type P, with the caution that nilpotence alone does not diagnose N. Step 3 adjoins a hyperbolic plane for type N, one generator for type C, and nothing for type P, the identity being the repair of an idempotent; a general type-P element has von Neumann pseudo-inverses, and among the idempotents with a given image the self-adjoint one is canonical exactly where a definite form is available. In the $(1,1)$ ascent the adjoined ladder pair $f, f^{*}$ and its occupation projectors are a system of matrix units, the adjoined pair is itself null, and the reflection in an adjoined direction is an automorphism of the ascent. \emph{Provenance:} \cite[Def.~3.1, Ex.~3.5, Thm.~3.9, Thm.~4.2, Rem.~5.8, Prop.~6.1, Rem.~6.6, Thm.~7.1, Prop.~7.8, Table~1]{ButlerT1}.
\end{assembled}
\begin{proof}
Every zero divisor of a nondegenerate real Clifford algebra is annihilated on each side by a nontrivial idempotent; in a simple algebra the idempotent is non-central and the obstruction is type P, refined to type N exactly when the annihilating idempotent is $wv$ for a hyperbolic pair $(v, w)$ of null vectors, $vw + wv = 1$. Step 2 of the algorithm diagnoses by three tests in order: a vector of square zero is type N; a central idempotent is type C; otherwise type P, with the caution that nilpotence alone does not diagnose N. Step 3 adjoins a hyperbolic plane for type N, one generator for type C, and nothing for type P, the identity being the repair of an idempotent; a general type-P element has von Neumann pseudo-inverses, and among the idempotents with a given image the self-adjoint one is canonical exactly where a definite form is available. In the $(1,1)$ ascent the adjoined ladder pair $f, f^{*}$ and its occupation projectors are a system of matrix units, the adjoined pair is itself null, and the reflection in an adjoined direction is an automorphism of the ascent. \emph{Provenance:} \cite[Def.~3.1, Ex.~3.5, Thm.~3.9, Thm.~4.2, Rem.~5.8, Prop.~6.1, Rem.~6.6, Thm.~7.1, Prop.~7.8, Table~1]{ButlerT1} (\ref{c:transport.no-transport-terminates}, \ref{c:transport.rule-1-morita-trivial}, \ref{c:transport.exhaustiveness-every-zero-divisor-annihilated-idempotent}, \ref{c:transport.type-p-self-repairing}).

This restates \ref{c:transport.transport-algorithm}.

\end{proof}

\begin{proposition}[The pairing creator under the transport algorithm]\label{c:descent-lift.pairing-creator-under-transport-algorithm}
Take two rungs in the root paper's real form, $\Cl(2,2)$ with orthogonal basis $e_{1}, e_{2}$ of square $+1$ and $e_{3}, e_{4}$ of square $-1$, ladder pairs $f_{k} = \tfrac12(e_{k} + e_{k+2})$, $f_{k}^{*} = \tfrac12(e_{k} - e_{k+2})$, number projectors $f_{k}^{*}f_{k}$ and vacuum projector $P_{00} = f_{1}f_{1}^{*}f_{2}f_{2}^{*}$ \cite[Thm.~7.1]{ButlerT1}. Let $n = f_{1}^{*}f_{2}^{*}$ be the pairing creator, $h = f_{1}^{*}f_{2}$ the hop, and $\dagger$ the Fock adjoint, the anti-automorphism that exchanges $f_{k}$ and $f_{k}^{*}$.
\begin{enumerate}
\item[(i)] $n$ is a nonzero pure bivector with $n^{2} = 0$, the algebra is simple, and $n$ has annihilating idempotents in the sense of Fact \ref{c:descent-lift.imported-taxonomy-type-p-rule-algorithm}, both non-central: $e = f_{2}f_{2}^{*}$ on the right, $n(1 - e) = n\,f_{2}^{*}f_{2} = 0$, and $f = f_{1}^{*}f_{1}$ on the left, $(1 - f)n = f_{1}f_{1}^{*}\,n = 0$. Under Step 2, which tests the element's grade first, $n$ is type P and Step 3 adjoins nothing. Both annihilating idempotents are of hyperbolic-pair form, $f_{2}f_{2}^{*} = wv$ for $(v, w) = (f_{2}^{*}, f_{2})$ and $f_{1}^{*}f_{1} = vw$ for $(v, w) = (f_{1}^{*}, f_{1})$, so under the refinement of Fact \ref{c:descent-lift.imported-taxonomy-type-p-rule-algorithm} the obstruction is the type-N one carried by the adjoined null vectors $f_{1}^{*}, f_{2}^{*}$ themselves, for which the no-go theorem closes every equivariant element-level repair and the point-level repair is a further $(1,1)$ ascent, which acts on points and not on $n$. On either reading the algorithm returns no element from $n$.
\item[(ii)] The pair annihilation $n^{\dagger} = f_{2}f_{1}$ satisfies the four Penrose equations with $n$: it is the Moore--Penrose inverse for $\dagger$, and its idempotents $nn^{\dagger} = P_{11}$ and $n^{\dagger}n = P_{00}$ are the self-adjoint ones that Fact \ref{c:descent-lift.imported-taxonomy-type-p-rule-algorithm} singles out where a definite form is available, the Fock form here. It annihilates the vacuum.
\item[(iii)] The bivector a rung carries, $f \wedge f^{*} = \tfrac12(ff^{*} - f^{*}f)$, squares to $+\tfrac14$ and commutes with the number operator: it is the Cartan element of the mode.
\item[(iv)] The reflection $e_{4} \mapsto -e_{4}$ extends to an automorphism of $\Cl(2,2)$ that fixes $f_{1}$, carries $f_{2}$ to $f_{2}^{*}$, $h$ to $n$, the hop generator $X_{h} = h - h^{\dagger}$ to the pairing generator $X_{p} = n - n^{\dagger}$, and $P_{00}$ to $P_{01}$.
\item[(v)] $X_{p}^{2} = -P_{\mathrm{even}}$ and $X_{h}^{2} = -P_{\mathrm{odd}}$ with $P_{\mathrm{even}} = P_{00} + P_{11}$ and $P_{\mathrm{odd}} = P_{01} + P_{10}$; so $\exp(\theta X) = (1 - P) + \cos\theta\,P + \sin\theta\,X$, and at $(\cos\theta, \sin\theta) = (\tfrac35, \tfrac45)$ both exponentials are rotors of $\mathrm{Spin}(2,2)$, unitary for $\dagger$; the pairing rotor carries $P_{00}$ to $\tfrac35 P_{00} + \tfrac45\,nP_{00}$, the singularity paper's witness state in the Fock ideal, and the hop rotor fixes $P_{00}$.
\item[(vi)] On the even sector $n$ and $n^{\dagger}$ satisfy the root paper's hyperbolic-pair identities, $n^{2} = (n^{\dagger})^{2} = 0$ and $nn^{\dagger} + n^{\dagger}n = P_{\mathrm{even}}$, and $X_{p} = n - n^{\dagger}$; the pairing rotation is the rotation of the plane this pair spans.
\item[(vii)] The root paper's pinching for either rung (Fact \ref{c:descent-lift.imported-pinching-one-way-door}) erases $n$, $h$, $X_{p}$ and $X_{h}$ alike and fixes the parity sections $\sigma_{1}, \sigma_{2}, \sigma_{1}\sigma_{2}$.
\end{enumerate}
\end{proposition}
\begin{proof}
Every clause is a finite identity in the blade basis of $\Cl(2,2)$, and the script checks each (Part H). For (ii) the adjoint is reversion composed with the automorphism $e_{3}, e_{4} \mapsto -e_{3}, -e_{4}$: an anti-automorphism, being an automorphism composed with one, and an involution, since the two commute; the script confirms the anti-automorphism property on random elements and that it exchanges $f_{k}$ with $f_{k}^{*}$. For (iv) the reflection is the automorphism the universal property attaches to the isometry, checked likewise on random elements. For (v), $X^{3} = -X$ and the idempotence of $P = -X^{2}$ are checked, and $XP = X$ follows, which gives the closed form; a rotor is an even element with $R\widetilde{R} = 1$. In the blade basis $X_{p} = \tfrac12(e_{1}e_{2} + e_{3}e_{4})$ and $X_{h} = \tfrac12(e_{1}e_{2} - e_{3}e_{4})$, checked: the two generators are the two combinations of the same two rotation planes, exchanged by reversing the orientation of one of them. The readings in (i) and (ii) rest on \ref{c:descent-lift.imported-taxonomy-type-p-rule-algorithm}, which supplies the type-P and type-N diagnosis, the repair Step 3 makes for each, and the canonical self-adjoint idempotents where a definite form is available; (vii) rests on \ref{c:descent-lift.imported-pinching-one-way-door}, which supplies the pinching $\EN(x) = PxP + (1-P)x(1-P)$ of each rung, whose kernel is the ladder sector.

\end{proof}

\phantomsection\label{c:note.descent-lift.setting-what-does-two-papers-what-3}
\textbf{What this settles.} Run on the pairing creator, the transport algorithm returns type P at cost zero and hands back the pair annihilation, an element of the singularity paper's vacuum stabiliser (Fact \ref{c:descent-lift.imported-descent-modes-univalence-supplied-class}); read through the type-N refinement instead, it returns the obstruction of the adjoined null vectors, which no equivariant element-level repair meets. Under Step 2 it adjoins nothing; under the refinement Step 3 adjoins a plane whose repair acts on points and returns no element from $n$; on neither reading is a rotation produced. The ascent itself contains the pairing and the hop on equal terms, as rotors of its own Spin group built from its own ladder elements, and it has a symmetry exchanging them (iv); that symmetry exchanges the parity sectors too, so it is a symmetry of the algebra and not an operation the univalence rule admits. So the transport principle neither produces nor excludes the pairing generator: what distinguishes a hop from a pairing is the vacuum, which of $f, f^{*}$ annihilates it, together with the parity grading the vacuum fixes, and the singularity paper's exclusion of pairing terms is a premise about that choice, not a consequence of transport. Clause (vi) says where the pairing generator lives: on the even sector of the pair, the creator and its pseudo-inverse are one effective ladder pair, Anderson's pseudospin \cite{Anderson1958}, and the pairing rotation is the rotation of that plane. The registered claim, that the principle produces the generator, is dead at Step 3 as the registration predicted; the kill conditions K1 to K3 of that registration triggered under Step 2, K4 held, every prediction P1 to P6 held, and clause (vii) is outside the registration.

\begin{remark}[Departures from the registered wording]\label{c:note.descent-lift.departures-registered-wording}
Five, recorded in one place. (1) The registration framed the count on the full Fock space, where the single-mode descent has Kraus rank $2$ and the $\nu$-mode descent rank $2^{\nu}$, and predicted a minimal record of $\nu + 1$ rungs (P1). On the physical, even states, which is where (P) places the framework and where the note is stated, the ranks are $1$ and $2^{\nu-1}$ and the minimal record is $\nu$ rungs; the single mode has nothing to lift, and the pair is the base case. The full-space count $\nu + 1$ is correct for the descent of unphysical coherences and is not the note's result. (2) K1 was framed for one mode, where on even states there is nothing to lift; its test, that the even record states must be orthogonal and one rung has one, applies unchanged to the pair, and Corollary \ref{c:descent-lift.parity-alone-forces-second-rung} records that both one-rung maps do induce the descent and fail by parity alone. (3) K4, P4, P5 and P6 were stated for a single mode $c\ket{0} + s\ket{1}$, which is not an even state; the note states them for the pair $c\ket{00} + s\ket{11}$, where every number predicted holds. (4) Claim D, K2 and P2 as registered place the pairing rotation on the record. Theorem \ref{c:descent-lift.lift-costs-one-rung-per-mode}(v) shows it may be placed on the system instead and carried to the record by hops; the clause that survives is that no number-conserving dynamics lifts, and the lift of Theorem \ref{c:descent-lift.lift-costs-one-rung-per-mode}(iv) places the pairing on the record as registered. (5) P7 as registered named the ledger entry transport-convenient; the grade found is transport-absent, for the reason given below.
\end{remark}

\phantomsection\label{c:note.descent-lift.setting-what-does-two-papers-what-4}
\textbf{Verdict.} K1 to K4 do not trigger, in the corrected framing: the two-rung count for the pair and the one-rung-per-mode count in general (Theorem \ref{c:descent-lift.lift-costs-one-rung-per-mode}(ii)); the necessity of a generator that fails to conserve number, a pairing among quadratic ones (Theorem \ref{c:descent-lift.lift-costs-one-rung-per-mode}(v)); the coupling in the supplied class (Theorem \ref{c:descent-lift.lift-costs-one-rung-per-mode}(iv)); the stored coherence with the pair's Schmidt coefficients (Corollary \ref{c:descent-lift.classical-shared-bit-reduced-view-pure}). K5 triggers, as registered: nothing in the root paper's transport principle forces a map to be lifted, so the step from ``the descent is one-way'' to ``the descent should be lifted'' is a reading of the principle, not a consequence of it. Of the predictions, P3 held as registered; P1, P4, P5 and P6 held with the framing corrected as Remark \ref{c:note.descent-lift.departures-registered-wording} records; P2 held as to the two quanta the record must gain and not as to where the pairing acts (item (4) of the remark); P7 held as to K5 and not as to the registered grade (item (5)).

\phantomsection\label{c:note.descent-lift.setting-what-does-two-papers-what-5}
\textbf{Prior art, and the two ledgers.} The dilation, the Kraus form and the minimal environment are Stinespring, Kraus and Choi; the parity constraint is the univalence rule in its operational form; the Thouless state and the Witt vacuum are imports. What the note proves for itself is the even-state Kraus rank $2^{\nu-1}$, the count of one rung per mode with the second rung forced by parity, the explicit lift and its factorisation into a supplied-class coupling and a pairing rotation, the identification of that rotation with the singularity paper's witness family, the number-conservation obstruction and its indifference to where the pairing acts, the stored-coherence and partial-lift computations, and Proposition \ref{c:descent-lift.within-supplied-dynamics-descent-has-no}. Under the two-ledger discipline the result is transport-absent: the rungs are where the record lives and the count matches the tower's one mode per rung, but the lift exists for any $\nu$ fermionic modes adjoined by any means, the univalence rule and Stinespring's theorem fix the count without the $(1,1)$ increment, and the transport principle does not force the lift (K5). The ascent supplies the vocabulary in which the record is counted, not a step the argument needs; and Proposition \ref{c:descent-lift.pairing-creator-under-transport-algorithm} shows that the creator from which the one generator the argument needs is built is one the principle classifies as costing nothing.

\phantomsection\label{c:note.descent-lift.setting-what-does-two-papers-what-6}
\textbf{What it does not do.} It does not touch the angle. The amplitudes $(c, s)$ of the pair are inputs to the lift, not outputs; the Hawking temperature enters through the singularity paper's regularity premise exactly as before. The record is an adjoined system, not a spatially separated one, and nothing here is a Bell test, a coupling Hamiltonian between separated regions, a global positive-frequency splitting, a non-Gaussian state, or an invariant measure of entanglement; the lift moves the pair's entanglement one cut outward and calibrates what that costs, and that is all it does.

\phantomsection\label{c:note.descent-lift.setting-what-does-two-papers-what-7}
\textbf{Consequences for the owning papers.} For the root paper: Theorem 7.6 can carry a clause, and Proposition \ref{c:descent-lift.pairing-creator-under-transport-algorithm}(i) records a point for it to settle: Step 2 diagnoses a product of two adjoined null vectors by its grade, as type P, while Theorem 3.9(2) refines it to type N by the origin of its annihilating idempotent; the two branches agree that no element is returned, and the root paper should say which diagnosis governs such products. The pinching's kernel, the ladder sector, is recoverable by an even isometry into $\nu$ further rungs, one per pinched mode, so the door is one-way for algebra homomorphisms and for the dynamics the completion supplies, and two-way for the tower if a pairing rotation is admitted. For the singularity paper: Premise 7.30's pairing clause now has a second characterisation beside Proposition 13.4's tangent: among quadratic generators the pairing terms are exactly those whose absence makes the descent irreversible, and Corollary 11.6's witness rotation is the record preparation of the descent's lift at amplitude $\tfrac43$, which overshoots (Corollary \ref{c:descent-lift.partial-lifts-witness-overshoots}); the complete lift is its member at amplitude $1$. Successor problems: whether the framework's own tower selects the amplitude $1$, that is, whether anything in the ascent prefers the maximally paired record; whether the lift composes along the tower, the record of the record costing $\nu$ further rungs each time; and the angle, which this note leaves where it found it.

\section{Position in the series}

\phantomsection\label{c:note.descent-lift.setting-position-series}
This document carries no T-number: it is a companion note, deposited with T5 \cite{ButlerT5} and read with T1 \cite{ButlerT1}, and its notation is that of the primer T0 \cite{ButlerT0}. From T1 it takes the pinching and the one-way door (Proposition 7.5, Theorem 7.6), the transport principle whose inversion it tries (Theorem 4.13), and, for Proposition \ref{c:descent-lift.pairing-creator-under-transport-algorithm}, the taxonomy (Definition 3.1, Theorem 3.9), the type-P rule and the pseudo-inverse remark (Proposition 6.1, Remark 6.6), the ladder pair as matrix units (Theorem 7.1) and the algorithm (Table 1). From T5 it takes the descent as the parity average and its multi-mode form (Theorem 5.7, Proposition 9.6), univalence (Premise 7.8, Fact 7.9), the supplied class (Premise 7.30, Theorem C.8), the stabiliser-tangent split (Proposition 13.4), the necessity witness (Corollary 11.6), and the thermofield pair with its sample point (Lemma 7.2, Theorem 7.27). From T3 \cite{ButlerT3} it takes one qubit per rung (Theorem 4.2), the count the lift reproduces. From T16 \cite{ButlerT16} it takes Cartan's theorem (Fact 2.10). It sends back to T1 a clause for Theorem 7.6 and to T5 a second characterisation of the pairing clause of Premise 7.30 and a reading of Corollary 11.6. The capstone T17 \cite{ButlerT17} carries the two-ledger grading used in the section ``What this does to the two papers, and what it does not do''.

\section{Machine verification}

\phantomsection\label{c:note.descent-lift.setting-machine-verification}
Every computational identity asserted in this note has been verified in exact arithmetic by the script \texttt{verify\_T05\_dilation.py} \cite{ButlerVerifyT05dilation}, which imports the series' shared library \texttt{GA.py}: its Gaussian rationals and its exact ranks. Irrational normalisations are avoided by working with unnormalised isometries, $W^{\dagger}W = 2^{\nu-1}$, and unnormalised record states; the reduced state of a subsystem is computed as the restriction of the joint state to the subsystem's algebra, with the operators $\ket{Y}\bra{X}_{S}$ built from the system's ladder operators, and the tensor-product partial trace is used only as a cross-check on even records. No floating point is used anywhere; the script exits nonzero on any failure. Thirty-four checks, all passing, in eight parts. The two pre-registrations are deposited beside the script \cite{ButlerPreregT05dilation, ButlerPreregT05pairing}.

\emph{Part A: the channel.} The product of the two parity averages equals the joint-number-diagonal pinching on all sixteen matrix units of the pair and is idempotent on a random matrix; on the even sector it equals $\tfrac12(\rho + \G_{0}\rho\G_{0}) = \tfrac12(\rho + \G_{1}\rho\G_{1})$, with Choi rank $2$ there and $4$ on the full space; the descent of the thermofield pair is $\operatorname{diag}(\tfrac{16}{25}, \tfrac{9}{25})$ with coherence $\tfrac{12}{25}$ removed; on one mode the even sector is the vacuum line and the descent is the identity on it.

\emph{Part B: the record.} The even sector of $k$ rungs has $2^{k-1}$ basis states for $k \leq 4$; Stinespring's one-rung map and the occupation copy both induce the descent on the even sector and both fail evenness, the first producing a parity-mixed joint state from the vacuum, the second an odd joint state from $\ket{11}$.

\emph{Part C: the lift of the pair.} $W = 1 \otimes \ket{00} + \G_{0} \otimes \ket{11}$ is even, has $W^{\dagger}W = 2$ with $W\ket{00} \perp W\ket{11}$, and induces $2\EN$ on the even sector's matrix units; the tensor-product partial trace agrees with the algebraic restriction on the lifted states; $C = 1 - 2P^{R}_{11}N_{0}$ is diagonal, equals $1 - \tfrac14(1-\G_{0})(1-\G_{2})(1-\G_{3})$, commutes with every $N_{k}$ and the total parity, squares to $1$; and $W = C(\,\cdot \otimes \varphi)$ on both even basis states and on the pair.

\emph{Part D: the preparation.} $\varphi$ is even and a Witt vacuum with Thouless amplitude $1$; the hops, the annihilation and the Cartan elements annihilate the record vacuum while the pair creation carries it to $\ket{11}$; the rational pairing rotor is unitary, even, and carries $\ket{00}$ to $\tfrac35\ket{00} + \tfrac45\ket{11}$; ten number-conserving generators commute with the total number, and the enumeration of joint basis states at the required number and parity leaves only the record vacuum for both even system states; a beam splitter between the pair and a rung at $(\tfrac35, \tfrac45)$ moves weight $\tfrac{16}{25}$ out of $\ket{11}$; the system pairing $1 + (a^{\dagger}_{0}a^{\dagger}_{1} - a_{1}a_{0})$, conjugated by the two system-record swaps and followed by $C$, is an even map inducing $2\EN$ on the even sector; the annihilation $a_{2}a_{3}$ kills the record vacuum while its Hermitian completion $a_{2}a_{3} + a^{\dagger}_{3}a^{\dagger}_{2}$ carries it to $\ket{11}$; the quartic element $iN_{0}(a^{\dagger}_{2}a^{\dagger}_{3} - a_{3}a_{2})$ is Hermitian, even, kills the joint vacuum and fails to commute with the total number, and its rotation at $\pi/2$, entered in closed form, is a unitary, even map carrying $\psi \otimes \ket{00}$ to the two-rung occupation copy, which is even, induces $\EN$, and equals $W$ in the record basis $\ket{00} \pm \ket{11}$.

\emph{Part E: the stored coherence.} The lifted pair is pure and even; its pair marginal is the classical shared bit and equals the descent; the marginal of each of the two modes is $\operatorname{diag}(\tfrac{16}{25}, \tfrac{9}{25})$; the record marginal has trace $1$ and determinant $\tfrac{144}{625}$; the lifted state equals $c\ket{00}\ket{+} + s\ket{11}\ket{-}$; the partial-lift factors $1, \tfrac{7}{25}, 0, -\tfrac{7}{25}$ at $A = 0, \tfrac34, 1, \tfrac43$.

\emph{Part F: the classical records.} The even mixture on the record pair with the coupling $C$ induces the descent exactly and leaves the joint output at purity $\tfrac12$; a single rung in the maximally mixed state with the coupling $1 - 2P^{R}_{1}N_{0}$ does the same, its joint output commuting with the total parity.

\emph{Part G: three modes.} On the even sector of three modes the pinching has Choi rank $4$; the even sector of three rungs has dimension $4$ and of two rungs $2$; the four record states are even and orthonormal (each pair checked); $W$ is even with $W^{\dagger}W = 4$ and induces $4\EN$ on the sixteen even matrix units; the record preparation is an even Witt vacuum with all three amplitudes $1$.

\emph{Part H: the pairing creator under the transport algorithm.} In $\Cl(2,2)$ over the rationals: $n = f_{1}^{*}f_{2}^{*}$ is nonzero, pure grade $2$, square zero, the centre is the scalars, $f_{2}^{*}f_{2}$ and $f_{1}^{*}f_{1}$ are non-central idempotents with $n\,f_{2}^{*}f_{2} = 0$ and $f_{1}f_{1}^{*}n = 0$, and the annihilating idempotents $f_{2}f_{2}^{*}$ and $f_{1}^{*}f_{1}$ are $wv$ and $vw$ for the grade-one null pairs $(f_{2}^{*}, f_{2})$ and $(f_{1}^{*}, f_{1})$ with $vw + wv = 1$; the Fock adjoint is an anti-automorphism exchanging $f_{k}$ and $f_{k}^{*}$, and $n^{\dagger} = f_{2}f_{1}$ satisfies the four Penrose equations with $n$ and kills $P_{00}$; $f \wedge f^{*}$ squares to $\tfrac14$ and commutes with the number operator; the reflection $e_{4} \mapsto -e_{4}$ is an automorphism carrying $h$ to $n$, $X_{h}$ to $X_{p}$, $P_{00}$ to $P_{01}$, $P_{\mathrm{odd}}$ to $P_{\mathrm{even}}$ and the hop rotor to the pairing rotor; $X_{p} = \tfrac12(e_{1}e_{2} + e_{3}e_{4})$ and $X_{h} = \tfrac12(e_{1}e_{2} - e_{3}e_{4})$, $X_{p}^{2} = -P_{\mathrm{even}}$, $X_{h}^{2} = -P_{\mathrm{odd}}$, $X^{3} = -X$, both closed-form exponentials at $(\tfrac35, \tfrac45)$ are even with $R\widetilde{R} = 1 = RR^{\dagger}$, the pairing rotor carries $P_{00}$ to $\tfrac35P_{00} + \tfrac45 nP_{00}$ and the hop rotor fixes it; $nn^{\dagger} + n^{\dagger}n = P_{\mathrm{even}}$ with $P_{\mathrm{even}} + P_{\mathrm{odd}} = 1$; and the pinching of either rung erases $n, h, X_{p}, X_{h}$ and fixes $\sigma_{1}, \sigma_{2}, \sigma_{1}\sigma_{2}$.

What the script does not check: Facts \ref{c:descent-lift.imported-pinching-one-way-door} to \ref{c:descent-lift.imported-thouless-states-witt-vacua} (imports); clauses (i) to (iii) of Theorem \ref{c:descent-lift.lift-costs-one-rung-per-mode} beyond $\nu = 2, 3$ and clauses (iv) and (v) beyond $\nu = 2$ (the proofs are by counting and by Fact \ref{c:descent-lift.imported-stinespring-kraus-choi}); the coefficient argument in the proof of Proposition \ref{c:descent-lift.within-supplied-dynamics-descent-has-no}, of which the script checks one annihilation and the one quartic exception; and, in Proposition \ref{c:descent-lift.pairing-creator-under-transport-algorithm}, the imported taxonomy of Fact \ref{c:descent-lift.imported-taxonomy-type-p-rule-algorithm} and the identity $XP = X$, which follows from $X^{3} = -X$.

\section{Results used from elsewhere}
The proofs above use the following results of the series, stated briefly; each is proved in the paper named.

\begin{fact}[Imported: Cartan's purity theorem and the Thouless form]\label{c:one-relation.imported-cartans-purity-theorem-thouless-form}
Imported from Cartan, Chevalley and Thouless: pure spinors are Witt vacua forming one orbit of dimension $\nu(\nu-1)/2$, all Weyl spinors are pure for $2\nu \leq 6$, and vacua overlapping $\ket{0}$ have the exponential Thouless form.
\emph{Source:} 'E. Cartan, Lecons sur la th'eorie des spineurs, Hermann (1938).; C. Chevalley, The Algebraic Theory of Spinors, Columbia University Press (1954).; D. J. Thouless, Stability conditions and nuclear rotations in the Hartree-Fock theory, Nucl. Phys. 21 (1960) 225--232..
\end{fact}

\begin{fact}[The class-relative no-go: no supplied generator hops]\label{c:singularity.class-relative-no-go-no-supplied}
Given function-linearity of the tame constructions, no generator in the supplied class moves occupation between points of $E$, while pairing generators lie outside the class and are not excluded.
\emph{Source:} T5, Thm C.8 (\href{https://doi.org/10.5281/zenodo.22721790}{doi:10.5281/zenodo.22721790}).
\end{fact}

\begin{fact}[The descent is the parity average]\label{c:singularity.descent-parity-average}
The descent is the parity average $E(x) = \tfrac12\bigl(x + \alpha(x)\bigr)$, the unique conditional expectation onto the even subalgebra, with kernel the odd ladder sector.
\emph{Source:} T5, Thm 5.7 (\href{https://doi.org/10.5281/zenodo.22721790}{doi:10.5281/zenodo.22721790}).
\end{fact}

\begin{fact}[The envelope and the endpoint, and why boundary-mode locality is needed for them]\label{c:singularity.envelope-endpoint-why-boundary-mode-locality}
Under (L), (U), (C), (I) and (L$_E$), $S_R(t) \leq \min(\Sth(t), C(t))$ with $S_R = 0$ at complete evaporation, and the bound can fail without (L$_E$).
\emph{Source:} T5, Cor 11.6 (\href{https://doi.org/10.5281/zenodo.22721790}{doi:10.5281/zenodo.22721790}).
\end{fact}

\begin{fact}[Imported: the univalence superselection rule]\label{c:singularity.imported-univalence-superselection-rule}
Imported: in a theory graded by fermion number, superpositions across total fermion parity are unobservable, so observables and physical states are even; the rule concerns total parity only.
\emph{Source:} G. C. Wick, A. S. Wightman and E. P. Wigner, The intrinsic parity of elementary particles, Phys. Rev. 88 (1952), 101--105.; O. Bratteli and D. W. Robinson, Operator Algebras and Quantum Statistical Mechanics 2: Equilibrium States. Models in Quantum Statistical Mechanics, 2nd ed., Springer (1997)., \S 5.2.
\end{fact}

\begin{fact}[(L$_E$) as the stabiliser and the tangent]\label{c:singularity.l-stabiliser-tangent}
In $\mathfrak{so}(2\nu)$ on Fock space, hops and pair annihilations stabilise the vacuum ray and pair creations span the tangent to its orbit, so only pairing can move the vacuum.
\emph{Source:} T5, Prop 13.4 (\href{https://doi.org/10.5281/zenodo.22721790}{doi:10.5281/zenodo.22721790}).
\end{fact}

\begin{fact}[(P$_{\mathrm{loc}}$), reduced: the metric-level reading on many modes]\label{c:singularity.p-reduced-metric-level-reading-many}
The mode-wise descents commute and compose to the pinching $\EN$, so if the metric level is read as all pointwise descents at once, it is the joint-number-diagonal algebra.
\emph{Source:} T5, Prop 9.6 (\href{https://doi.org/10.5281/zenodo.22721790}{doi:10.5281/zenodo.22721790}).
\end{fact}

\begin{fact}[Purity is an output: the partner is the commutant]\label{c:singularity.purity-output-partner-commutant}
For a faithful single-mode state the GNS construction supplies the partner as the commutant and, for a Gibbs state, the thermofield double, so purity and the partner are outputs of the one-sided state.
\emph{Source:} T5, Thm 7.27 (\href{https://doi.org/10.5281/zenodo.22721790}{doi:10.5281/zenodo.22721790}).
\end{fact}

\begin{fact}[Schmidt forcing, proved outright]\label{c:singularity.schmidt-forcing-proved-outright}
A unit vector on the doubled mode that is boost invariant up to phase and cyclic and separating is $\cos\theta\,|00\rangle + \sin\theta\,|11\rangle$ with both coefficients nonzero, so its Schmidt basis is the number basis.
\emph{Source:} T5, Lem 7.2 (\href{https://doi.org/10.5281/zenodo.22721790}{doi:10.5281/zenodo.22721790}).
\end{fact}

\begin{fact}[The door is one-way]\label{c:transport.door-one-way}
There is no unital algebra homomorphism from either minimal completion of a nondegenerate $\Cl(p,q)$ back to $\Cl(p,q)$ (the Rule 2 one existing only when $p - q \equiv 1 \pmod 4$), so every descent is non-multiplicative.
\emph{Source:} T1, Thm 7.6 (\href{https://doi.org/10.5281/zenodo.21743863}{doi:10.5281/zenodo.21743863}).
\end{fact}

\begin{fact}[Exhaustiveness: every zero divisor is annihilated by an idempotent]\label{c:transport.exhaustiveness-every-zero-divisor-annihilated-idempotent}
Every nonzero zero divisor of a nondegenerate real Clifford algebra is annihilated on each side by a nontrivial idempotent, so the N, P, C taxonomy of obstructions is exhaustive.
\emph{Source:} T1, Thm 3.9 (\href{https://doi.org/10.5281/zenodo.21743863}{doi:10.5281/zenodo.21743863}).
\end{fact}

\begin{fact}[No transport terminates]\label{c:transport.no-transport-terminates}
Every algebra reached by at least one transport step has zero divisors, of type N after Rule 1 and of type P after Rule 2, so transport never reaches a division algebra.
\emph{Source:} T1, Prop 7.8 (\href{https://doi.org/10.5281/zenodo.21743863}{doi:10.5281/zenodo.21743863}).
\end{fact}

\begin{fact}[The pinching, and its kernel]\label{c:transport.pinching-kernel}
The pinching by the descent projector is a conditional expectation onto its commutant whose kernel is exactly the adjoined ladder sector, and it returns transported elements unchanged.
\emph{Source:} T1, Prop 7.5 (\href{https://doi.org/10.5281/zenodo.21743863}{doi:10.5281/zenodo.21743863}).
\end{fact}

\begin{fact}[Rule 1 is Morita-trivial]\label{c:transport.rule-1-morita-trivial}
Rule 1 gives $\Cl(p+1, q+1) \cong \Mat_2(\Cl(p,q))$ with the adjoined ladder pair as off-diagonal matrix units, so it preserves the Morita class, the Wedderburn type and the clock residue $p - q \pmod 8$.
\emph{Source:} T1, Thm 7.1 (\href{https://doi.org/10.5281/zenodo.21743863}{doi:10.5281/zenodo.21743863}).
\end{fact}

\begin{fact}[The transport algorithm]\label{c:transport.transport-algorithm}
The transport results assemble into a six-step procedure (locate, diagnose, apply the rule, compute, descend, iterate) in which every step is licensed by a proved result and no case is left uncovered.
\emph{Source:} T1, Sec 8.1 (\href{https://doi.org/10.5281/zenodo.21743863}{doi:10.5281/zenodo.21743863}).
\end{fact}

\begin{fact}[Type P is self-repairing]\label{c:transport.type-p-self-repairing}
In any unital ring the identity assignment $\sigma(e) = e$ on idempotents satisfies both von Neumann identities and is equivariant under every automorphism of every extension preserving the ring, so type P obstructions are repaired at signature increment $(0,0)$.
\emph{Source:} T1, Prop 6.1 (\href{https://doi.org/10.5281/zenodo.21743863}{doi:10.5281/zenodo.21743863}).
\end{fact}

\begin{thebibliography}{99}
\bibitem{ButlerT1} L.~Butler,
\emph{The Transport Principle: Total Division in Clifford Algebras via Minimal Completion},
Transport Papers T1 (2026).
\bibitem{ButlerT5} L.~Butler,
\emph{Division by Zero at the Gravitational Singularity: The Transport Principle Applied to the Black Hole, and the Information Paradox as a Category Error},
Transport Papers T5 (2026), doi:\href{https://doi.org/10.5281/zenodo.22721790}{10.5281/zenodo.22721790}.
\bibitem{ButlerT3} L.~Butler,
\emph{The Tower and the CAR Limit: Iterated Transport, the Frozen Clock, and Second Quantisation},
Transport Papers T3 (2026), doi:\href{https://doi.org/10.5281/zenodo.22721556}{10.5281/zenodo.22721556}.
\bibitem{ButlerT17} L.~Butler,
\emph{Capstone: The Transport Programme, Method, Status, and the Falsifiable Ledger},
Transport Papers T17 (2026).
\bibitem{WWW1952} G.~C.~Wick, A.~S.~Wightman and E.~P.~Wigner,
The intrinsic parity of elementary particles,
\emph{Phys.\ Rev.} \textbf{88} (1952), 101--105.
\bibitem{Stinespring1955} W.~F. Stinespring,
\emph{Positive functions on $C^{*}$-algebras},
Proc. Amer. Math. Soc. 6 (1955) 211--216.
\bibitem{Kraus1983} K.~Kraus,
\emph{States, Effects, and Operations: Fundamental Notions of Quantum Theory},
Lecture Notes in Physics 190, Springer (1983).
\bibitem{Choi1975} M.-D. Choi,
\emph{Completely positive linear maps on complex matrices},
Linear Algebra Appl. 10 (1975) 285--290.
\bibitem{BravyiKitaev2002} S.~B. Bravyi, A.~Yu. Kitaev,
\emph{Fermionic quantum computation},
Ann. Phys. 298 (2002) 210--226.
\bibitem{ArakiMoriya2003} H.~Araki, H.~Moriya,
\emph{Joint extension of states of subsystems for a CAR system},
Commun. Math. Phys. 237 (2003) 105--122.
\bibitem{Thouless1960} D.~J. Thouless,
\emph{Stability conditions and nuclear rotations in the Hartree-Fock theory},
Nucl. Phys. 21 (1960) 225--232.
\bibitem{ButlerT16} L.~Butler,
\emph{Quantum Mechanics and General Relativity from One Relation: The Two Shadows and the Laboratory},
Transport Papers T16 (2026), doi:\href{https://doi.org/10.5281/zenodo.22723212}{10.5281/zenodo.22723212}.
\bibitem{ButlerPreregT05dilation} L.~Butler,
\emph{preregistration\_T05\_dilation.md: the kill conditions and predictions of ``The Descent, Lifted'', as registered before computation},
deposited with this note (2026).
\bibitem{ButlerPreregT05pairing} L.~Butler,
\emph{preregistration\_T05\_pairing\_type.md: the kill conditions and predictions of the pairing-creator experiment of ``The Descent, Lifted'', as registered before computation},
deposited with this note (2026).
\bibitem{Anderson1958} P.~W. Anderson,
\emph{Random-phase approximation in the theory of superconductivity},
Phys. Rev. 112 (1958) 1900--1916.
\bibitem{ButlerT0} L.~Butler,
\emph{An Introduction to Geometric Algebra: A Primer for the Transport Papers},
Transport Papers T0 (2026).
\bibitem{ButlerVerifyT05dilation} L.~Butler,
\emph{verify\_T05\_dilation.py: exact-arithmetic verification of the computational identities of ``The Descent, Lifted''},
Python script, deposited with this note alongside the shared library GA.py (2026).
\end{thebibliography}
\end{document}